\documentclass[10pt,a4paper]{amsart}
\usepackage{amsmath,amssymb,amsthm,graphicx}
\usepackage{verbatim}
\newenvironment{reference}{\comment}{\endcomment}
%\newenvironment{reference}{}{}
\newenvironment{slogan}{\comment}{\endcomment}
\newenvironment{history}{\comment}{\endcomment}

% For commutative diagrams we use Xy-pic
\usepackage[all]{xy}

% We use 2cell for 2-commutative diagrams.
\xyoption{2cell}
% Theorem environments.
%
\theoremstyle{plain}
\newtheorem{theorem}[subsection]{Theorem}
\newtheorem{proposition}[subsection]{Proposition}
\newtheorem{lemma}[subsection]{Lemma}

\theoremstyle{definition}
\newtheorem{definition}[subsection]{Definition}
\newtheorem{example}[subsection]{Example}
\newtheorem{exercise}[subsection]{Exercise}
\newtheorem{situation}[subsection]{Situation}

\theoremstyle{remark}
\newtheorem{remark}[subsection]{Remark}
\newtheorem{remarks}[subsection]{Remarks}

\numberwithin{equation}{subsection}

% Macros
%
\def\lim{\mathop{\mathrm{lim}}\nolimits}
\def\colim{\mathop{\mathrm{colim}}\nolimits}
\def\Spec{\mathop{\mathrm{Spec}}}
\def\Hom{\mathop{\mathrm{Hom}}\nolimits}
\def\Ext{\mathop{\mathrm{Ext}}\nolimits}
\def\SheafHom{\mathop{\mathcal{H}\!\mathit{om}}\nolimits}
\def\SheafExt{\mathop{\mathcal{E}\!\mathit{xt}}\nolimits}
\def\Sch{\mathit{Sch}}
\def\Mor{\mathop{\mathrm{Mor}}\nolimits}
\def\Ob{\mathop{\mathrm{Ob}}\nolimits}
\def\Sh{\mathop{\mathit{Sh}}\nolimits}
\def\NL{\mathop{N\!L}\nolimits}
\def\CH{\mathop{\mathrm{CH}}\nolimits}
\def\proetale{{pro\text{-}\acute{e}tale}}
\def\etale{{\acute{e}tale}}
\def\QCoh{\mathit{QCoh}}
\def\Ker{\mathop{\mathrm{Ker}}}
\def\Im{\mathop{\mathrm{Im}}}
\def\Coker{\mathop{\mathrm{Coker}}}
\def\Coim{\mathop{\mathrm{Coim}}}

% Boxtimes
%
\DeclareMathSymbol{\boxtimes}{\mathbin}{AMSa}{"02}

%
% Macros for moduli stacks/spaces
%
\def\QCohstack{\mathcal{QC}\!\mathit{oh}}
\def\Cohstack{\mathcal{C}\!\mathit{oh}}
\def\Spacesstack{\mathcal{S}\!\mathit{paces}}
\def\Quotfunctor{\mathrm{Quot}}
\def\Hilbfunctor{\mathrm{Hilb}}
\def\Curvesstack{\mathcal{C}\!\mathit{urves}}
\def\Polarizedstack{\mathcal{P}\!\mathit{olarized}}
\def\Complexesstack{\mathcal{C}\!\mathit{omplexes}}
% \Pic is the operator that assigns to X its picard group, usage \Pic(X)
% \Picardstack_{X/B} denotes the Picard stack of X over B
% \Picardfunctor_{X/B} denotes the Picard functor of X over B
\def\Pic{\mathop{\mathrm{Pic}}\nolimits}
\def\Picardstack{\mathcal{P}\!\mathit{ic}}
\def\Picardfunctor{\mathrm{Pic}}
\def\Deformationcategory{\mathcal{D}\!\mathit{ef}}
\usepackage{cleveref}
\crefname{originalreference}{original reference}{original references}
\Crefname{originalreference}{Original reference}{Original references}
\makeatletter
\newlabel{example-degree-simple-algebraic-extension}{{09G7}{source}{}{sourcepointer-1}{}}
\newlabel{example-degree-simple-algebraic-extension@cref}{{[originalreference][0][]09G7}{[1][0][]source}{}{}{}}
\makeatother
\title{Four equivalent characterizations of algebraically closed fields}
\author{The Stacks Project Authors}
\date{stacks-project:fields.tex:local-snapshot-20261012; GNU FDL 1.2 or later}
\begin{document}
\maketitle
\noindent\textit{Editorial note.} All delivered mathematical text is the author's own, copied from the pinned Stacks Project chapter file fields.tex; no statement or proof step was rewritten, reordered or removed.\par\medskip
\noindent\textit{Editorial note.} The counted proof is the author's complete argument: The chapter's own definition-algebraically-closed (source lines 1117-1121, triviality of every algebraic extension) is characterization (1); the author notes before the lemma that other texts use the root formulation instead.\par\medskip
\noindent\textit{Editorial note.} The label lemma-algebraically-closed names the counted statement, and the proof environment immediately following it in the source is delivered verbatim as its proof.\par\medskip
\noindent\textit{Editorial note.} The counted statement is carried by the author's own lemma environment in the chapter source.\par\medskip
\noindent\textit{Allowed standard prerequisites (editorial).}\begin{itemize}
\item Same chapter: example-degree-simple-algebraic-extension (a root of an irreducible polynomial of degree d gives the extension k[x]/(P) of degree d, so an irreducible of degree > 1 yields a nontrivial finite algebraic extension); division with remainder in F[x]
\end{itemize}\medskip
\ifdefined\hypertarget\hypertarget{sourcepointer-1}{}\fi
\noindent\textit{Source reference.} Original reference 09G7: sources/stacks-fields/pinned-original-source.tex, source lines 721--721. Prerequisite pointer invoked at 1142; it names a standard fact of the same chapter that the counted argument uses, and every step of the counted proof built on it is retained verbatim in the delivered body.\par\medskip
\begin{lemma}
\label{lemma-algebraically-closed}
Let $F$ be a field. The following are equivalent
\begin{enumerate}
\item $F$ is algebraically closed,
\item every irreducible polynomial over $F$ is linear,
\item every nonconstant polynomial over $F$ has a root,
\item every nonconstant polynomial over $F$ is a product of linear factors.
\end{enumerate}
\end{lemma}

\begin{proof}
If $F$ is algebraically closed, then every irreducible polynomial is linear.
Namely, if there exists an irreducible polynomial of degree $> 1$, then
this generates a nontrivial finite (hence algebraic) field extension, see
Example \ref{example-degree-simple-algebraic-extension}.
Thus (1) implies (2). If every irreducible polynomial
is linear, then every irreducible polynomial has a root, whence every
nonconstant polynomial has a root. Thus (2) implies (3).

\medskip\noindent
Assume every nonconstant polynomial has a root. Let $P \in F[x]$
be nonconstant. If $P(\alpha) = 0$ with $\alpha \in F$, then we see
that $P = (x - \alpha)Q$ for some $Q \in F[x]$ (by division with remainder).
Thus we can argue by induction on the degree that any nonconstant
polynomial can be written as a product $c \prod (x - \alpha_i)$.

\medskip\noindent
Finally, suppose that every nonconstant polynomial over $F$ is a product of
linear factors. Let $E/F$ be an algebraic extension. Then all the simple
subextensions $F(\alpha)/F$ of $E$ are necessarily trivial (because the
only irreducible polynomials are linear by assumption). Thus $E = F$.
We see that (4) implies (1) and we are done.
\end{proof}

\end{document}
